\documentclass{article}
\usepackage[utf8]{inputenc}
\usepackage{amsmath,amssymb}
\usepackage[most]{tcolorbox}
\usepackage{geometry}
\geometry{margin=2.5cm}

\newcommand{\step}[1]{\par\noindent\hangindent=3.5em\hangafter=1%
  \makebox[3.5em][l]{\textbf{#1.}}}
\newcommand{\steptag}[1]{\hfill{\small\textsf{[#1]}}}

\newtcolorbox{proofbox}[1]{
  colback=gray!5, colframe=gray!60,
  fonttitle=\bfseries, title={#1},
  breakable, enhanced,
  left=4pt, right=4pt, top=4pt, bottom=4pt,
  before skip=10pt, after skip=10pt
}

\begin{document}

\begin{proofbox}{Uniform compressed-unit action: there are universal C\_co ...}
\step{1} Uniform compressed-unit action: there are universal C\_co < infinity and e\_co > 0 such that, for e=delta+epsilon <= e\_co, every compatible amplified rectangular corner satisfies ||u\_T dot A-A|| <= C\_co*e*||A|| and ||A dot u\_R-A|| <= C\_co*e*||A||. \steptag{validated} \\[6pt]
\step{1.1} Setup and uniform primitive bounds: interpret a compatible amplified rectangular corner as an extended epsilon-C*-algebra calA, n>=1, delta-projections T,R, T\_n=I\_n tensor T, R\_n=I\_n tensor R, and A in S\_\{T\_n,R\_n\}; write u\_\{T,n\}=Co\_\{T\_n,T\_n\}(T\_n) and u\_\{R,n\}=Co\_\{R\_n,R\_n\}(R\_n). There are universal M,k<infinity and e\_0>0 such that, when e=delta+epsilon<=e\_0, ||T\_n||,||R\_n||<=M, both compression maps Co\_\{T\_n,R\_n\} are within k*e of L\_\{T\_n\}R\_\{R\_n\} and R\_\{R\_n\}L\_\{T\_n\}, and A=Co\_\{T\_n,R\_n\}(A). \steptag{validated} \\[6pt]
\step{1.1.1} Amplified admissibility and notation. By def-extended-epsilon-cstar-algebra, M\_n tensor calA is an epsilon-C*-algebra. The block diagonals T\_n=I\_n tensor T and R\_n=I\_n tensor R are delta-projections because adjoints and products amplify entrywise and ||I\_n tensor Z||=||Z||. By lem-compcb-amplified-compression, the relevant maps and spaces are the amplified compressions and corners; by lem-compcb-amplified-compression-identities each amplified Co is idempotent. Hence A in S\_\{T\_n,R\_n\}=Im(Co\_\{T\_n,R\_n\}) implies A=Co\_\{T\_n,R\_n\}(A), while u\_\{T,n\}=Co\_\{T\_n,T\_n\}(T\_n) lies in S\_\{T\_n,T\_n\}, and similarly for R. \steptag{validated} \\[4pt]
\step{1.1.2} Uniform constants. Unpack the universal O-terms in def-delta-projection and def-compressed-corner: there are universal a,k<infinity and e\_0>0 such that every delta-projection S in an epsilon-C*-algebra with e<=e\_0 obeys either ||S||<=a*delta or ||S||<=1+a*e, and ||Co\_\{S,V\}-L\_S R\_V||,||Co\_\{S,V\}-R\_V L\_S||<=k*e for every compatible delta-projection V. Shrink e\_0<=1 and set M=max(a,1+a). Applying this inside M\_n tensor calA gives ||T\_n||,||R\_n||<=M and all assertions of node 1.1. \steptag{validated} \\[4pt]
\step{1.2} Uniform compressed-unit comparison: after possibly decreasing e\_0, there is a universal K\_u<infinity such that ||u\_\{T,n\}-T\_n||<=K\_u*e, ||u\_\{R,n\}-R\_n||<=K\_u*e, and ||u\_\{T,n\}||,||u\_\{R,n\}||<=M+K\_u for every compatible amplified rectangular corner. \steptag{validated} \\[6pt]
\step{1.2.1} Single-projection estimate. Let S be either T\_n or R\_n and u\_S=Co\_\{S,S\}(S). From node 1.1, ||S||<=M and ||Co\_\{S,S\}-L\_S R\_S||<=k*e. Since (L\_S R\_S)(S)=S(S\textasciicircum{}2), the epsilon-banach-cstar-norm-axioms and ||S\textasciicircum{}2-S||<=delta give ||u\_S-S||<=k*e*||S||+||S(S\textasciicircum{}2)-S||, while bilinearity gives S(S\textasciicircum{}2)-S=S(S\textasciicircum{}2-S)+(S\textasciicircum{}2-S), so ||u\_S-S||<=[k*M+2*M+1]*e when e<=1. \steptag{validated} \\[4pt]
\step{1.2.2} Taking K\_u=k*M+2*M+1 in node 1.2 and using the triangle inequality yields ||u\_\{T,n\}||,||u\_\{R,n\}||<=M+K\_u*e<=M+K\_u. The constant and threshold are universal and independent of n. \steptag{validated} \\[6pt]
\step{1.2.2.1} Validated node 1.1.2 supplies universal finite constants M and k, independent of n, and a universal threshold e\_0<=1 with ||T\_n||,||R\_n||<=M. Validated node 1.2.1, applied first with S=T\_n and then with S=R\_n, supplies ||u\_\{T,n\}-T\_n||<=(k*M+2*M+1)*e and ||u\_\{R,n\}-R\_n||<=(k*M+2*M+1)*e. Define K\_u=k*M+2*M+1. For S=T\_n and S=R\_n separately, the triangle inequality gives ||u\_S||<=||u\_S-S||+||S||<=K\_u*e+M<=M+K\_u because 0<=e<=e\_0<=1. Hence both bounds asserted in node 1.2.2 hold with a universal finite K\_u and a universal threshold, independently of n; no premise from the pending parent node 1.2 is used. \steptag{validated} \\[4pt]
\step{1.3} Ambient rectangular unit action: there is a universal K\_a<infinity such that, under the setup of node 1.1 and e<=e\_0, ||u\_\{T,n\}A-A||<=K\_a*e*||A|| and ||A u\_\{R,n\}-A||<=K\_a*e*||A||. \steptag{validated} \\[6pt]
\step{1.3.1} Left projection action. Put D=T\_n(A R\_n). From A=Co\_\{T\_n,R\_n\}(A) and ||Co\_\{T\_n,R\_n\}-L\_\{T\_n\}R\_\{R\_n\}||<=k*e in node 1.1, ||A-D||<=k*e||A||. The identity T\_n A-A=T\_n(A-D)+(T\_nD-D)+(D-A), the epsilon-banach-cstar-norm-axioms, and ||T\_n||,||R\_n||<=M give ||T\_n(A-D)||<=2*M*k*e||A|| and ||A R\_n||<=2*M||A||. Writing Y=A R\_n, approximate associativity and ||T\_n\textasciicircum{}2-T\_n||<=delta give ||T\_nD-D||=||T\_n(T\_nY)-T\_nY||<=[epsilon*M\textasciicircum{}2+2*delta]*||Y||<=2*M*(M\textasciicircum{}2+2)*e||A||. Thus ||T\_n A-A||<=K\_T*e||A|| for universal K\_T=2*M*k+2*M*(M\textasciicircum{}2+2)+k. \steptag{validated} \\[4pt]
\step{1.3.2} Right projection action. Put D=(T\_n A)R\_n. The second compression comparison in node 1.1 gives ||A-D||<=k*e||A||. Now A R\_n-A=(A-D)R\_n+(D R\_n-D)+(D-A). Using ||T\_n A||<=2*M||A||, approximate associativity to compare ((T\_nA)R\_n)R\_n with (T\_nA)(R\_n\textasciicircum{}2), and ||R\_n\textasciicircum{}2-R\_n||<=delta yields the same kind of universal estimate ||A R\_n-A||<=K\_R*e||A||, with for instance K\_R=2*M*k+2*M*(M\textasciicircum{}2+2)+k. \steptag{validated} \\[4pt]
\step{1.3.3} Replace projections by compressed units. By node 1.2 and the multiplication bound in epsilon-banach-cstar-norm-axioms, ||(u\_\{T,n\}-T\_n)A||<=2*K\_u*e||A|| and ||A(u\_\{R,n\}-R\_n)||<=2*K\_u*e||A||. Combining these with nodes 1.3.1 and 1.3.2 proves node 1.3 with K\_a=max(K\_T,K\_R)+2*K\_u. \steptag{validated} \\[4pt]
\step{1.3.4} Dependency-closed compressed-unit comparison, independent of node 1.2. Let S be T\_n or R\_n and put u\_S=Co\_\{S,S\}(S). Validated node 1.1 gives ||S||<=M, ||S\textasciicircum{}2-S||<=delta, and e<=e\_0<=1; the universal O(e) compression comparison in registered def-compressed-corner gives a universal k<infinity (enlarging the setup constant if necessary) with ||Co\_\{S,S\}-L\_S R\_S||<=k*e. Since (L\_S R\_S)(S)=S(S\textasciicircum{}2), the operator-norm estimate and the epsilon-banach-cstar multiplication bound give ||u\_S-S||<=k*e*||S||+||S(S\textasciicircum{}2)-S||. By bilinearity, S(S\textasciicircum{}2)-S=S(S\textasciicircum{}2-S)+(S\textasciicircum{}2-S), so ||S(S\textasciicircum{}2)-S||<=(1+epsilon)||S||*||S\textasciicircum{}2-S||+||S\textasciicircum{}2-S||<=(2*M+1)*e. Hence, with the universal K\_u\_local=k*M+2*M+1, both ||u\_\{T,n\}-T\_n||<=K\_u\_local*e and ||u\_\{R,n\}-R\_n||<=K\_u\_local*e follow using only validated node 1.1 and registered definitions. \steptag{validated} \\[4pt]
\step{1.3.5} Dependency-closed assembly. By node 1.3.4 and e<=1, ||(u\_\{T,n\}-T\_n)A||<=(1+epsilon)||u\_\{T,n\}-T\_n||||A||<=2*K\_u\_local*e||A||, and likewise ||A(u\_\{R,n\}-R\_n)||<=2*K\_u\_local*e||A||. Bilinearity yields u\_\{T,n\}A-A=(u\_\{T,n\}-T\_n)A+(T\_nA-A) and A u\_\{R,n\}-A=A(u\_\{R,n\}-R\_n)+(A R\_n-A). Combining the preceding estimates with validated nodes 1.3.1 and 1.3.2 proves both bounds in node 1.3 with the universal K\_a=max(K\_T,K\_R)+2*K\_u\_local. This derivation does not invoke pending node 1.2. \steptag{validated} \\[4pt]
\step{1.4} Compressed-product assembly: combining the ambient estimates with lem-compcb-rectangular-product yields universal C\_co<infinity and e\_co>0 for which ||u\_\{T,n\} dot A-A||<=C\_co*e*||A|| and ||A dot u\_\{R,n\}-A||<=C\_co*e*||A||; these u\_\{T,n\},u\_\{R,n\} are the amplified compressed units denoted u\_T,u\_R in the root. \steptag{validated} \\[6pt]
\step{1.4.1} The pairs (u\_\{T,n\},A) in S\_\{T\_n,T\_n\} x S\_\{T\_n,R\_n\} and (A,u\_\{R,n\}) in S\_\{T\_n,R\_n\} x S\_\{R\_n,R\_n\} are compatible amplified rectangular pairs by node 1.1. Therefore lem-compcb-rectangular-product supplies universal C\_p<infinity and e\_p>0 such that, whenever e=delta+epsilon<=e\_p, ||u\_\{T,n\} dot A-u\_\{T,n\}A||<=C\_p*e*||u\_\{T,n\}||||A|| and ||A dot u\_\{R,n\}-A u\_\{R,n\}||<=C\_p*e*||A||||u\_\{R,n\}||. \steptag{validated} \\[4pt]
\step{1.4.2} Choose e\_co as the minimum of e\_0, the thresholds in lem-compcb-amplified-compression, lem-compcb-amplified-compression-identities, and lem-compcb-rectangular-product, and 1; choose C\_co=K\_a+C\_p*(M+K\_u). Nodes 1.2-1.4.1 and the triangle inequality then give both root bounds. All constants and e\_co are universal and independent of n, calA, T,R,A; by lem-compcb-amplified-compression, u\_\{T,n\}=Co\_\{T\_n\}(T\_n)=I\_n tensor Co\_T(T) and similarly for R, so this is exactly the root notation. \steptag{validated} \\[6pt]
\step{1.4.2.1} Dependency-closed assembly via children 1.4.2.1.1 and 1.4.2.1.2. The first child derives uniform bounds ||u\_\{T,n\}||,||u\_\{R,n\}||<=M+K\_u and ambient action bounds ||u\_\{T,n\}A-A||,||A u\_\{R,n\}-A||<=K\_a*e*||A|| directly from validated nodes 1.2.1, 1.3.1, and 1.3.2, without using pending nodes 1.2 or 1.3. The second child combines those bounds with validated node 1.4.1. Explicitly, for e<=e\_co:=min(e\_0,e\_p,e\_cmp,e\_id,1), the triangle inequality gives ||u\_\{T,n\} dot A-A||<=[C\_p(M+K\_u)+K\_a]e||A|| and ||A dot u\_\{R,n\}-A||<=[C\_p(M+K\_u)+K\_a]e||A||. Thus C\_co:=K\_a+C\_p(M+K\_u) works. The constants and thresholds are universal, and lem-compcb-amplified-compression identifies the amplified compressed units with the root notation. \steptag{validated} \\[6pt]
\step{1.4.2.1.1} Local bounds from validated inputs only. Put K\_u:=k*M+2*M+1. Validated node 1.2.1, applied to S=T\_n and S=R\_n, gives ||u\_\{T,n\}-T\_n||<=K\_u*e and ||u\_\{R,n\}-R\_n||<=K\_u*e whenever e<=e\_0<=1. Together with ||T\_n||,||R\_n||<=M from the validated setup, the triangle inequality gives ||u\_\{T,n\}||,||u\_\{R,n\}||<=M+K\_u. The multiplication bound ||XY||<=(1+epsilon)||X||||Y|| and epsilon<=e<=1 give ||(u\_\{T,n\}-T\_n)A||<=2*K\_u*e*||A|| and ||A(u\_\{R,n\}-R\_n)||<=2*K\_u*e*||A||. Bilinearity yields u\_\{T,n\}A-A=(u\_\{T,n\}-T\_n)A+(T\_nA-A) and A u\_\{R,n\}-A=A(u\_\{R,n\}-R\_n)+(A R\_n-A). Hence validated nodes 1.3.1 and 1.3.2 imply ||u\_\{T,n\}A-A||<=K\_a*e*||A|| and ||A u\_\{R,n\}-A||<=K\_a*e*||A|| for K\_a:=max(K\_T,K\_R)+2*K\_u. Thus all four estimates are obtained without invoking pending nodes 1.2 or 1.3. \steptag{validated} \\[4pt]
\step{1.4.2.1.2} Final assembly from the preceding local bounds and validated node 1.4.1. Let e\_co:=min(e\_0,e\_p,e\_cmp,e\_id,1), where e\_p is the positive threshold in node 1.4.1 and e\_cmp,e\_id are the positive thresholds in lem-compcb-amplified-compression and lem-compcb-amplified-compression-identities. If e<=e\_co, node 1.4.1 and the preceding child give ||u\_\{T,n\} dot A-A||<=||u\_\{T,n\} dot A-u\_\{T,n\}A||+||u\_\{T,n\}A-A||<=[C\_p(M+K\_u)+K\_a]e||A||. Likewise ||A dot u\_\{R,n\}-A||<=||A dot u\_\{R,n\}-A u\_\{R,n\}||+||A u\_\{R,n\}-A||<=[C\_p(M+K\_u)+K\_a]e||A||. Therefore C\_co:=K\_a+C\_p(M+K\_u) works. Every constant is finite and universal and every threshold in the minimum is positive and universal, so C\_co and e\_co are universal and independent of n,calA,T,R,A. Finally lem-compcb-amplified-compression identifies u\_\{T,n\}=Co\_\{T\_n\}(T\_n)=I\_n tensor Co\_T(T), and similarly for R, which is the root notation. \steptag{validated} \\[4pt]
\end{proofbox}

\end{document}
